\documentclass[12pt, titlepage]{article}
%\usepackage[norsk]{babel}
\usepackage[utf8]{inputenc}
\usepackage{minted, enumerate}
\usepackage[colorlinks = true,
linkcolor = blue,
urlcolor = blue]{hyperref}


\begin{document}

\title{Refleksjonsnotat}
\author{IIRA2001 - Høst 2023}
\date{}
\maketitle

I løpet av semesteret har vi sett på grunnleggende programmering i python, og anvendelse av dette i dataanalyse og økonomiske fag.
Siste del av mappen består av et \textit{refleksjonsnotat} hvor dere reflekterer rundt læringsprosessen dere har vært gjennom, og hvordan programmering kan være faglig relevant i videre studie og virke.

Et refleksjonsnotat er en personlig tekst hvor du reflekterer over egne tanker, erfaringer eller innsikter knyttet til en bestemt hendelse eller prosess (slik som læringsprosessen dere nå har vært gjennom).
Ved å bearbeide og reflektere rundt disse tankene og erfaringene er målet å fremme læring og personlig og faglig vekst.

~\\
I hovedtrekk skal dere reflektere rundt hvordan programmering kan være faglig relevant, hvilke ferdiheterer dere har tilegnet dere og hvordan dere ser for dere fremtidig utvikling og bruk av det dere har lært.\\
%\textbf{Tips, relevante punkter å ha med, og relevant spørsmål å reflektere rundt:}
\subsubsection*{Tips, relevante punkter og spørsmål}
\begin{itemize}
   \item{Legg vekt på refleksjon, analytisk og kritisk tenkning}
   \item{Pek ut de viktigste programmeringsferdighetene du har utviklet}
   \item{Hvordan har programmering styrket forståelsen av økonomiske eller andre faglige aspekter?}
   \item{Har noen av oppgavene eller casene bidratt forbedret forståelse av økonomiske problemstillinger?}
   \item{Hvilke utfordringer møtte du, og hvordan håndterte du disse?}
   \item{Har du lært noe som er relevant for en karriere du ser for deg?}
   \item{Du kan diskutere potensielle anvendelser}
   \item{Identifiser noen ting som var vellykkede, og hvorfor}
   \item{Hvordan kan du bygge på det du har lært i fremtiden?}
   \item{Er det kanskje noe i oppgavene du har levert du kunne ønske å forbedre?}
   \item{Diskuter og grei ut om hva du har lært}
   \item{Du kan reflektere over egen tilnærming til å lære programmering}
   \item{Det vil være lurt å knytte refleksjoner til konkrete eksempler i mappen}
   \item{Refleksjonsnotat er en personlig tekst, så vær ærlig og åpen}
   \item{Det kan være lurt med en oppsummering og en slags konklusjon på slutten}
\end{itemize}

Du trenger naturligvis ikke å oppfylle alle punktene over, men de kan være til hjelp for å komme i gang.

~\\
\textit{Refleksjonsnotatet burde være på rundt 750-1500 ord}

\end{document}
